\chapter{Male mastectomy}
\index{Male mastectomy}

\section*{Definition and Overview}
Male mastectomy is the surgical removal of breast tissue in a male patient, primarily performed to treat male breast cancer or severe, refractory gynecomastia (benign enlargement of male breast tissue). The procedure involves removing the breast gland, including the breast parenchyma and fatty tissue, and is classified as simple or total mastectomy (removing breast tissue and the nipple-areola complex), subcutaneous mastectomy (sparing the nipple-areola complex), or modified radical mastectomy (removing breast tissue, the nipple-areola complex, and regional axillary lymph nodes). While breast cancer in males is relatively rare, its surgical management mirrors that of female breast cancer. Subcutaneous mastectomy is also a key surgical option for gender-affirming care in transgender men. The surgery is critical to eradicate localised malignancy, prevent systemic metastasis, and restore chest contour, significantly improving clinical and psychological outcomes.

\section*{Epidemiology}
Male breast cancer is an uncommon malignancy, accounting for approximately 1\% of all breast cancer diagnoses globally. In many countries, approximately 150 to 200 men are diagnosed with breast cancer annually, typically presenting at an older age than women, with a peak incidence between 60 and 70 years. Because of a lack of screening programmes and lower public awareness, men often present at a later stage, necessitating mastectomy rather than breast-conserving surgery. Gynecomastia, by contrast, is highly prevalent, affecting up to 30\% to 60\% of adolescent boys and older men; however, only a small percentage of patients with severe, painful, or medically refractory cases undergo surgical resection. The epidemiology of male mastectomy is also shaped by rising rates of gender-affirming chest reconstruction (mastectomy), reflecting increased access to transgender healthcare.

\section*{Aetiology and Pathophysiology}
The indications for male mastectomy stem from neoplastic genetic pathways, hormonal imbalances, or gender dysphoria.
\begin{itemize}
    \item \textbf{Genetic mutations:} Highly correlated with BRCA2 mutations (and to a lesser extent BRCA1), which dramatically increase the lifetime risk of male breast cancer to approximately 6--8\%.
    \item \textbf{Hormonal imbalance:} An elevated oestrogen-to-androgen ratio, which can be caused by liver disease (cirrhosis), obesity, Klinefelter syndrome (XXY), or oestrogen-producing testicular tumours.
    \item \textbf{Pharmacological triggers:} Long-term use of medications that induce gynecomastia, such as spironolactone, anabolic steroids, anti-androgens for prostate cancer, or certain antiretroviral therapies.
    \item \textbf{Malignant transformation:} Clonal proliferation of ductal epithelial cells, leading most commonly to invasive ductal carcinoma (IDC) in males.
    \item \textbf{Gender congruence:} The physiological mismatch between an individual's gender identity and their anatomical chest characteristics, necessitating chest masculinisation surgery.
\end{itemize}

\section*{Clinical Features}
The clinical presentation of patients undergoing male mastectomy depends on the underlying pathology.
\begin{itemize}
    \item \textbf{Painless subareolar mass:} The most common presentation of male breast cancer, typically presenting as a firm, fixed, unilateral mass eccentric to the nipple.
    \item \textbf{Nipple retraction or discharge:} Inward pulling of the nipple or spontaneous serous or bloody discharge, indicative of underlying ductal pathology.
    \item \textbf{Skin changes:} Dimpling, puckering, or ulceration of the skin overlying the breast tissue, or erythema resembling the skin of an orange (peau d'orange).
    \item \textbf{Gynecomastia symptoms:} Unilateral or bilateral, concentric, rubbery or firm tissue beneath the nipple-areola complex, often accompanied by tenderness during the active proliferative phase.
    \item \textbf{Axillary lymphadenopathy:} Palpable, firm nodes in the ipsilateral axilla, indicating regional metastatic spread of breast cancer.
\end{itemize}

\section*{Differential Diagnosis}
It is essential to differentiate male breast cancer from benign breast conditions and other chest wall masses.
\begin{itemize}
    \item \textbf{Gynecomastia:} Benign proliferation of glandular tissue, which is typically soft or rubbery, symmetric, and situated directly behind the nipple-areola complex.
    \item \textbf{Pseudogynecomastia:} Deposition of fat in the breast area without glandular tissue proliferation, commonly seen in obese males.
    \item \textbf{Lipoma:} A slow-growing, benign fatty tumour that is mobile and soft, situated in the subcutaneous tissue.
    \item \textbf{Sebaceous or epidermal cyst:} A benign cutaneous fluid-filled sac, often with a central punctum, superficial to the breast parenchyma.
    \item \textbf{Chest wall tumour:} Primary skeletal or soft-tissue neoplasms (e.g. sarcoma) originating from the ribs or intercostal muscles.
\end{itemize}

\section*{When to Seek Medical Care}
Any male patient who notices a new, firm, or changing breast mass, nipple changes, or unexplained discharge should seek prompt evaluation by a general practitioner or breast specialist. Immediate emergency services should be contacted for severe postoperative complications.

\textbf{Contact your local emergency services for:}
\begin{itemize}
    \item \textbf{Severe postoperative haemorrhage:} Rapidly expanding, tense swelling of the chest wall (haematoma) accompanied by skin discolouration, active bleeding, or shortness of breath.
    \item \textbf{Signs of systemic sepsis:} High fever, severe chills, rapid heart rate, or confusion accompanying chest wound inflammation.
    \item \textbf{Anaphylaxis:} Severe allergic reactions to postoperative analgesics or antibiotics, presenting with difficulty breathing or facial swelling.
\end{itemize}

\section*{Diagnosis and Investigations}
A triple assessment approach is utilised to diagnose male breast conditions, combining clinical, radiological, and pathological findings.
\begin{itemize}
    \item \textbf{Clinical breast examination:} Detailed bilateral palpation of the breast tissue, nipple-areola complex, axillary, and supraclavicular lymph node regions.
    \item \textbf{Mammography and ultrasound:} High-resolution imaging to distinguish glandular tissue from fat (gynecomastia vs. pseudogynecomastia) and identify suspicious microcalcifications or solid masses.
    \item \textbf{Core needle biopsy:} The definitive diagnostic tool, obtaining tissue samples for histopathology, receptor status (ER, PR, HER2), and grading of suspected malignancies.
    \item \textbf{Genetic testing:} Referral for BRCA1 and BRCA2 genetic counselling and testing, recommended for all men diagnosed with breast cancer.
    \item \textbf{Staging scans (CT, bone scan, or PET-CT):} Indicated for diagnosed male breast cancer to evaluate for distant pulmonary, hepatic, or skeletal metastases.
\end{itemize}

\section*{Management}
Management involves surgical resection combined with oncological or supportive therapies, depending on the surgical indication.
\begin{itemize}
    \item \textbf{Mastectomy surgery:} Total or simple mastectomy for early-stage cancer; subcutaneous mastectomy (sparing the nipple-areola complex) for gynecomastia or gender-affirming surgery.
    \item \textbf{Axillary staging:} Sentinel lymph node biopsy (SLNB) or axillary lymph node clearance (ALNC) performed concurrently for invasive breast cancer.
    \item \textbf{Adjuvant systemic therapy:} Endocrine therapy (typically Tamoxifen) for hormone receptor-positive cancers, chemotherapy, or targeted HER2 therapies based on tumour biology.
    \item \textbf{Adjuvant radiotherapy:} Recommended following mastectomy for patients with large tumours, positive surgical margins, or multiple involved lymph nodes.
    \item \textbf{Postoperative drain management:} Monitoring and management of temporary closed-suction drains (e.g. Jackson-Pratt) placed to collect fluid and prevent seroma formation.
\end{itemize}

\section*{Complications}
\begin{itemize}
    \item \textbf{Seroma or haematoma formation:} Accumulation of fluid or blood under the skin flaps, often requiring needle aspiration or surgical evacuation.
    \item \textbf{Wound dehiscence or flap necrosis:} Separation of wound edges or death of the skin flaps due to compromised local blood supply.
    \item \textbf{Lymphedema:} Persistent swelling of the ipsilateral arm, occurring primarily in patients who undergo axillary lymph node clearance.
    \item \textbf{Nipple-areola necrosis:} Loss of viability of the nipple-areola complex, a risk in subcutaneous or nipple-sparing mastectomies.
    \item \textbf{Chronic chest wall pain:} Persistent neuralgia or phantom breast pain due to transection of intercostal nerves during tissue dissection.
\end{itemize}

\section*{Prognosis and Outcomes}
The prognosis for male mastectomy varies depending on the indication. For gynecomastia and gender-affirming surgery, outcomes are excellent, with high rates of patient satisfaction and significant improvements in body image and quality of life. For male breast cancer, the prognosis is highly dependent on the stage at diagnosis. The 5-year survival rate for localised male breast cancer is approximately 90\%, which drops significantly if lymph nodes are involved or distant metastasis has occurred. Because males have very little breast tissue, tumours quickly invade the overlying skin or pectoralis major muscle, highlighting the critical importance of early detection and prompt mastectomy to optimise survival.

\section*{Prevention and Self-Care}
\begin{itemize}
    \item \textbf{Perform regular self-checks:} Men, especially those with a family history of breast cancer or known BRCA mutations, should regularly check their chest for lumps.
    \item \textbf{Manage body weight:} Maintaining a healthy body weight to reduce oestrogen levels and the risk of gynecomastia and breast cancer.
    \item \textbf{Limit alcohol consumption:} Restricting alcohol intake to support liver health, preventing gynecomastia and hormonal imbalances.
    \item \textbf{Strict postoperative wound care:} Keeping the incision site clean and dry, and wearing compression garments as directed.
    \item \textbf{Gentle range-of-motion exercises:} Performing prescribed shoulder and arm stretches post-drain removal to prevent joint stiffness.
\end{itemize}

\section*{Sources and Further Reading}
The following reputable organisations provide further information on male mastectomy and male breast health:
\begin{itemize}
    \item Breast Cancer Network Australia. \textit{Support networks and resources specifically for men diagnosed with breast cancer.} https://www.bcna.org.au/
    \item Cancer Council Australia. \textit{Male breast cancer guides, risk factors, and treatment.} https://www.cancer.org.au/
    \item Mayo Clinic. \textit{Mastectomy procedures, risks, and recovery.} https://www.mayoclinic.org/
    \item American Cancer Society. \textit{Breast cancer in men: treatment options and surgery details.} https://www.cancer.org/
    \item Healthdirect Australia. \textit{Gynecomastia causes, symptoms, and surgical management.} https://www.healthdirect.gov.au/
\end{itemize}
